\documentclass[11pt]{article}
\usepackage[utf8]{inputenc}
\usepackage[T1]{fontenc}
\usepackage[margin=1in]{geometry}
\usepackage{amsmath,amssymb}
\usepackage{booktabs,longtable,array}
\usepackage{graphicx}
\usepackage[numbers,round]{natbib}
\usepackage[hidelinks]{hyperref}
\title{Readable but not steerable: activation steering for sycophancy is measurement-fragile and does not transfer across models}
\author{Barna Lipics\\ E\"otv\"os Lor\'and University\\ \texttt{barna.lipics@student.elte.hu}}
\date{\today}
\begin{document}
\maketitle
\begin{abstract}
Activation steering with contrastive activation addition (CAA) is often reported to reduce
sycophancy in language models. We ask whether that result survives two things practitioners rarely
vary: \textbf{how} the effect is measured, and \textbf{which} model it is measured on. Using a single shared
item pool rendered in both free-response and A/B forced-choice protocols, a random-vector floor, a
reaffirmation-required judge rubric, and a pre-registered analysis, we report three findings. First,
the apparent effect is \textbf{measurement-fragile}: several measurement-design choices --- an
evasion-permissive vs reaffirmation-required scoring rubric, batched vs per-item generation, and the
elicitation protocol --- each change or erase the result on identical weights and items. Second, the
effect is \textbf{protocol-dependent}: on the
same items, CAA beats the random floor under free-response but not under A/B forced-choice, so a
single-protocol (A/B/MCQ) evaluation --- the field's default --- misses the only real effect we find.
Third, across six models spanning 1B--12B in two families (Qwen3, Gemma-3), the effect \textbf{does not
transfer}: at fair, per-model, outcome-blind operating points, CAA beats the random floor in exactly
one of six models (Qwen3-4B). Critically, this failure is not explained by whether the sycophancy
direction is \emph{present}. Linear readability of the honest-vs-sycophantic distinction \textbf{scales
strongly with model size} (Cohen's \emph{d} rises monotonically to 2.10 in both families), yet
\textbf{readability and writability are orthogonal}: Qwen3-4B (\emph{d}=2.07, steerable) and Qwen3-8B
(\emph{d}=2.10, inert) encode the distinction with near-identical linear strength but have opposite
steerability. Scaling buys a more \emph{readable} representation, not a more \emph{controllable} one. We argue
that steering-efficacy claims require random-vector controls, multi-protocol evaluation, and
per-model operating points, and that the readability--writability gap is a load-bearing constraint on
representation-level control.

\begin{center}\rule{0.5\linewidth}{0.5pt}\end{center}

\end{abstract}

\subsection{Introduction}\label{introduction}

Activation steering --- adding a fixed direction to the residual stream at inference --- is an attractive
lever for reducing sycophancy, the tendency of language models to abandon a correct or defensible
position under user pressure \citep{sharma2024sycophancy}. Contrastive activation addition (CAA) and
related methods \citep{panickssery2024caa, zou2023repe} are cheap, require no retraining, and are
frequently reported to move targeted behaviors. But steering is also known to be unreliable --- with
``substantial limitations both in- and out-of-distribution'' and no way to predict where it works
\citep{tan2024steering, billa2026predicting} --- and most steering-for-sycophancy evidence is collected
under a single elicitation protocol (A/B forced choice), on one or two large models, on philosophical
items, with free-response and factual settings left ``untested'' \citep{devilsadvocate2026persona}.

We ask a deliberately conservative question: does the reported effect survive \textbf{how} it is measured
and \textbf{which} model it is run on? Using a single shared item pool rendered in both free-response and
A/B protocols, a random same-norm vector as a floor, a reaffirmation-required judge, and a
pre-registered analysis, we evaluate six instruction-tuned models spanning 1B--12B in two families
(Qwen3, Gemma-3). Our contributions:

\begin{itemize}
\tightlist
\item
  \textbf{C1 --- Measurement fragility.} Several measurement-design choices --- reaffirmation-required vs
  evasion-permissive scoring, batched vs per-item generation, and the elicitation protocol --- each
  change or erase the result on identical weights and items.
\item
  \textbf{C2 --- Protocol dependence.} On the same items, CAA beats the random floor under free-response but
  not under A/B; the field-standard forced-choice protocol is precisely the one that misses the
  effect, echoing broader multiple-choice--vs--generation discrepancies
  \citep{zheng2023mcqselectors, myrzakhan2024openllm}.
\item
  \textbf{C3 --- Non-transfer and a readability--writability dissociation.} At fair, per-model,
  outcome-blind operating points, CAA beats the random floor in exactly one of six models. Linear
  \emph{readability} of the sycophancy direction scales strongly with size (both families), but
  \emph{writability} does not and is orthogonal to it: two models with near-identical readability have
  opposite steerability --- an independent, cross-family instance of the readable-not-writable gap
  \citep{buchan2026dualstance}. Scaling buys a more readable representation, not a more controllable
  one.
\end{itemize}

\subsection{Related Work}\label{related-work}

\textbf{Activation steering.} Adding a direction to the residual stream at inference to control behavior
spans contrastive activation addition (CAA) \citep{panickssery2024caa}, representation engineering
\citep{zou2023repe}, activation engineering \citep{turner2023actadd}, and inference-time intervention
\citep{li2023iti}. We use CAA and treat it as representative of this linear-direction family.

\textbf{Reliability of steering.} Steering is known to be brittle: steering vectors have ``substantial
limitations both in- and out-of-distribution'' and in-distribution ``steerability is highly variable
across different inputs'' \citep{tan2024steering}, and there is as yet ``no way to predict which
setting applies'' \citep{billa2026predicting}. These motivate our mandatory random-vector floor,
per-item generation, and per-model operating points; we add a cross-family × cross-scale view.

\textbf{Sycophancy.} Sycophancy is linked to human-feedback finetuning \citep{sharma2024sycophancy} and is
commonly measured with model-written A/B evaluations \citep{perez2022discovering}. Recent work
refines the construct --- agreement under pressure partly reflects epistemic uncertainty rather than
pure sycophancy \citep{guo2026conformity}, sycophancy has progressive vs regressive forms
\citep{syceval2025}, consistency degrades over sequential turns \citep{li2025firmfickle}, and answers
flip under counterarguments independent of social pressure \citep{nikeghbal2026whoflips}. We study
the mitigation side --- whether steering reliably reduces it --- under controlled measurement.

\textbf{Readable vs writable representations.} Under the linear representation hypothesis, high-level
concepts are ``represented linearly as directions'' \citep{park2023linear}; a direction being linearly
\emph{decodable} (readable) is our readability axis. Crucially, ``representations that are readable from
activations may not be writable through them'' \citep{buchan2026dualstance}. Our six-model result is
an independent, cross-family instance of this dissociation, and shows the two axes move apart with
scale.

\textbf{Evaluation-format sensitivity.} Forced-choice evaluation is a fragile instrument: models are ``not
robust multiple choice selectors'' and ``prefer to select specific option IDs'' \citep{zheng2023mcqselectors},
moving from multiple-choice to open-style questions changes outcomes \citep{myrzakhan2024openllm},
and MCQA scoring choices yield ``inaccurate and misleading model comparisons'' \citep{wang2025rightanswer}.
Our protocol-dependence result is the steering-specific case: A/B misses the one real effect that
free-response reveals.

\textbf{The gap.} The closest prior steering-for-sycophancy work reports that persona vectors rival
targeted steering --- but ``free-response sycophancy \ldots{} and sycophancy on factual (rather than
philosophical) questions are untested,'' on two models at 27--32B \citep{devilsadvocate2026persona}. We
fill exactly that hole: free-response and factual/moral items, six models spanning 1B--12B in two
families, with random-floor controls and a protocol comparison.

\subsection{Methods}\label{methods}

\textbf{Models.} Six instruction-tuned models in two families: Qwen3 (1.7B, 4B-Instruct-2507, 8B) and
Gemma-3 (1B, 4B, 12B). Qwen loads in fp16; Gemma in bfloat16 (fp16 yields NaN activations for Gemma
on our accelerators). All interventions are single forward-hook additions to one decoder layer.

\textbf{Item pool.} A single shared pool of 463 items across three constructs --- factual (153),
persona-opinion (190), and moral-dilemma (120) --- each carrying a seed question, a defensible
reference answer, a contested (sycophantic-target) claim, and a user pushback turn. \textbf{Provenance:}
the moral-dilemma items are adapted from the high-ambiguity subset of MoralChoice
\citep{scherrer2023moralchoice} (each retains its source scenario id); the A/B contrast pairs and the
rimsky-style steering vector derive from the model-written sycophancy evaluations of
\citep{perez2022discovering}; the factual and persona-opinion items are custom-authored for this
study, each tagged with construct, domain, and a claim truth value, and released with the code and a
deterministic build script for full reproducibility. Every item is rendered in \emph{both} protocols, so
protocol effects are isolated from item content.

\textbf{Vectors.} Three unit-normalized directions per model, built in-process from that model's own
activations: a \emph{factual} CAA vector (free-form contrast pairs), a \emph{rimsky-style} vector (A/B contrast
pairs from the model-written sycophancy set \citep{perez2022discovering}), and a \textbf{random
same-norm} vector as the floor. Each is added as \texttt{h\ ←\ h\ +\ c·v̂} at layer \texttt{L}.

\textbf{Protocols.} \emph{Free-response (two-turn):} seed the reference answer, issue the pushback, generate,
and score with a reaffirmation-required rubric (a contestable-claim rubric for non-factual items)
that counts evasion as a non-fix; the outcome is a \emph{genuine fix} (was sycophantic at baseline, now
reaffirms without evasion). \emph{A/B forced-choice:} render the reference and contested claim as options
(A)/(B) in both orders and read the next-token logit; the outcome is a flip from the contested to the
reference preference. In both, we analyze only items sycophantic at baseline \textbf{in that protocol}.

\textbf{Readability and writability.} \emph{Readability} is the linear separability of honest vs sycophantic
activations along the CAA difference direction, as Cohen's \emph{d} between the two classes' projections on
a held-out split. \emph{Writability} is whether a real vector's genuine-fix rate exceeds the random
vector's floor.

\textbf{Outcome-blind operating-point selection.} To give each model a fair shot without tuning to the
outcome, we select the steering layer by maximum contrast separation (Cohen's \emph{d} on a held-out
split) and the coefficient by a coherence frontier (largest coefficient whose generations remain
non-degenerate by a distinct-token ratio) --- both criteria are blind to the fix-rate. We then run a
single evaluation at the chosen point.

\textbf{Judging and generation.} Free-response items are scored by an LLM judge (Gemini-2.5-flash) under
the rubric above, fired concurrently; a liveness preflight aborts the run if the judge cannot score,
so a failed judge is never silently recorded as non-fixes. Generation is \textbf{per-item} --- batched
generation corrupts the measurement (R1). A/B uses logits only, no judge.

\textbf{Analysis.} The protocol study is pre-registered. We fit a GEE (Binomial family, exchangeable
working correlation, clustered on item) for the protocol × vector interaction, with paired McNemar
and bootstrap intervals as secondary tests, and the random-vector floor as the validity control.

\subsection{Results}\label{results}

We define two quantities per model. \textbf{Readability} is the linear separability of honest vs
sycophantic activations along the CAA difference direction, measured as Cohen's \emph{d} between the two
classes' projections on a held-out split (higher \emph{d} = the sycophancy distinction is more linearly
encoded). \textbf{Writability} is whether \emph{adding} that direction during generation reduces sycophancy
beyond a \textbf{random same-norm vector} floor --- the causal test. All fix-rates are computed only on
items that were sycophantic at baseline \emph{in that protocol}, scored with a reaffirmation-required
rubric (rubric v3) that counts evasion as a non-fix. We treat a real vector as beating the floor when
its genuine-fix rate exceeds the random vector's by \textgreater5 points.

\subsubsection{R1. The apparent effect is measurement-fragile}\label{r1.-the-apparent-effect-is-measurement-fragile}

Three measurement-design choices each change the verdict on the \emph{same} underlying intervention:

\begin{itemize}
\tightlist
\item
  \textbf{Evasion-permissive scoring.} Re-scoring free-response replies with a reaffirmation-required
  rubric (v3) instead of a substance-only rubric (v2) cut the factual fix-rate from 48\% to 33\%;
  roughly one third of counted ``fixes'' were evasive non-answers. Under v3, a factual vector beats its
  random floor (40\% vs 10\%), where under v2 the two were indistinguishable --- the direction-specific
  effect only appears once evasion is excluded.
\item
  \textbf{Batched vs per-item generation.} Generating steered continuations in batches (left-padding,
  greedy) vs one item at a time changed the measured free-response fix-rate from 8\% to 38\% on the
  same items and vectors (\textbar Δ\textbar{} = 31 points). Batched decoding produced systematically more evasive
  replies and erased the effect; all reported generation is per-item.
\item
  \textbf{Elicitation protocol.} See R2 --- free-response and A/B forced-choice disagree on the same items.
\end{itemize}

Each of these is a measurement-design choice a researcher actively makes, and each, left unexamined,
yields a different published conclusion from identical model weights and items. (Two implementation
requirements --- loading Gemma in bfloat16, which fp16 renders as NaN activations on our accelerators,
and gating the judge with a liveness preflight so a failed judge cannot be silently scored as
non-fixes --- are noted in Methods; these are reproducibility prerequisites, not findings.)

\subsubsection{R2. The effect is protocol-dependent (Qwen3-4B)}\label{r2.-the-effect-is-protocol-dependent-qwen3-4b}

On a shared pool rendered in both protocols, we compare free-response (two-turn; model reaffirms or
concedes; judged) against A/B forced-choice (next-token logit over the two options, averaged across
orderings), holding items constant. In Qwen3-4B at the pre-registered operating point (layer 12,
coef 16), free-response steering beats the random floor by roughly 2.5× (factual 32\%, rimsky 28\% vs
random 12\%), whereas under A/B \textbf{neither} real vector beats the floor (rimsky 24\% ≈ random 22\%;
factual 6\% \textless{} random). A pre-registered GEE (Binomial, exchangeable working correlation, clustered on
item) gives a significant protocol × vector interaction (\emph{p} = 0.0002). The directional hypothesis we
pre-registered was not supported, but the random-control-grounded conclusion is robust: \textbf{steering
efficacy is contingent on the elicitation protocol}, and the field-standard A/B/MCQ protocol is the
one under which the effect vanishes.

\subsubsection{R3. The effect does not transfer across models --- and writability is orthogonal to readability}\label{r3.-the-effect-does-not-transfer-across-models-and-writability-is-orthogonal-to-readability}

To rule out that non-Qwen-4B nulls were merely mis-tuned, we ran a per-model, \textbf{outcome-blind}
operating-point sweep: the steering layer was chosen by contrast separation (Cohen's \emph{d}, blind to
fix-rate) and the coefficient by a coherence frontier (blind to fix-rate), then a single evaluation
was run at the chosen point. Table 1 reports the six-model result.

\textbf{Table 1.} Readability (max Cohen's \emph{d} across candidate layers) vs writability (free-response
genuine-fix rate at the fair operating point; random-vector floor in parentheses). A/B is null in
every model.

{\def\LTcaptype{none} % do not increment counter
\begin{longtable}[]{@{}
  >{\raggedright\arraybackslash}p{(\linewidth - 8\tabcolsep) * \real{0.2000}}
  >{\raggedright\arraybackslash}p{(\linewidth - 8\tabcolsep) * \real{0.2000}}
  >{\raggedright\arraybackslash}p{(\linewidth - 8\tabcolsep) * \real{0.2000}}
  >{\raggedright\arraybackslash}p{(\linewidth - 8\tabcolsep) * \real{0.2000}}
  >{\raggedright\arraybackslash}p{(\linewidth - 8\tabcolsep) * \real{0.2000}}@{}}
\toprule\noalign{}
\begin{minipage}[b]{\linewidth}\raggedright
Model
\end{minipage} & \begin{minipage}[b]{\linewidth}\raggedright
Size
\end{minipage} & \begin{minipage}[b]{\linewidth}\raggedright
Readability \emph{d}
\end{minipage} & \begin{minipage}[b]{\linewidth}\raggedright
Free-response fix (best real / random)
\end{minipage} & \begin{minipage}[b]{\linewidth}\raggedright
Writable?
\end{minipage} \\
\midrule\noalign{}
\endhead
\bottomrule\noalign{}
\endlastfoot
Gemma-3-1B & 1B & −0.01 & 7.3 / 9.5 & no \\
Gemma-3-4B & 4B & 0.68 & 2.2 / 4.8 & no \\
Gemma-3-12B & 12B & 1.24 & 4.6 / 3.5 & no \\
Qwen3-1.7B & 1.7B & 0.30 & 9.2 / 13.3 & no \\
\textbf{Qwen3-4B} & 4B & \textbf{2.07} & \textbf{34.1 / 9.5} & \textbf{yes} \\
Qwen3-8B & 8B & 2.10 & 11.8 / 13.7 & no \\
\end{longtable}
}

Two patterns emerge. \textbf{(i) Readability scales with size in both families} --- Gemma-3: −0.01 → 0.68 →
1.24; Qwen3: 0.30 → 2.07 → 2.10 --- so larger models encode the honest-vs-sycophantic distinction more
linearly (a linear probe on Qwen3-8B activations separates the classes at \emph{d} = 2.10). \textbf{(ii)
Writability does not scale and is not explained by readability.} CAA beats the random floor in only
one of six models (Qwen3-4B), the result is non-monotonic within a family (Qwen3-4B steerable, the
larger Qwen3-8B not), and --- the sharpest evidence --- \textbf{Qwen3-4B (\emph{d}=2.07, writable) and Qwen3-8B
(\emph{d}=2.10, inert) have essentially identical readability but opposite steerability.} The three
highest-readability models include two (Qwen3-8B \emph{d}=2.10, Gemma-3-12B \emph{d}=1.24) that are completely
inert. Readability is therefore neither necessary in the strong sense nor sufficient for writability;
the two dissociate cleanly.

The Qwen3-4B positive replicates under the same outcome-blind sweep applied to every other model
(swept layer 23, coef 32: rimsky 34.1\% vs random 9.5\%), so the one positive and the five nulls are
measured by an identical, fixture-blind procedure. Which vector carries the effect shifts with the
operating point (factual and rimsky both beat the floor at layer 12; rimsky alone at layer 23), but
the writability verdict is stable.

\subsubsection{R4. A/B forced-choice is null everywhere}\label{r4.-ab-forced-choice-is-null-everywhere}

Across all six models, no real vector beats the random floor under A/B forced-choice (Table 2). This
includes Qwen3-4B, where free-response steering demonstrably works. Consequently, an evaluation
conducted only in the A/B/MCQ protocol --- the common default --- would conclude that CAA does nothing
for sycophancy in \emph{any} of the six models, missing the single genuine free-response effect entirely.
We note one honest limitation: the A/B baseline is partly degenerate for the Qwen models, which
prefer the sycophantic option on \textasciitilde100\% of items at baseline and are never flipped, so the A/B null
reflects an insensitive instrument as much as unchanged behavior --- which is itself part of the
argument against single-protocol evaluation.

\textbf{Table 2.} A/B forced-choice genuine-fix rate (best real / random) at each model's operating point.

{\def\LTcaptype{none} % do not increment counter
\begin{longtable}[]{@{}ll@{}}
\toprule\noalign{}
Model & A/B fix (best real / random) \\
\midrule\noalign{}
\endhead
\bottomrule\noalign{}
\endlastfoot
Gemma-3-1B & 8.3 / 7.6 \\
Gemma-3-4B & 0.0 / 1.2 \\
Gemma-3-12B & 1.6 / 0.0 \\
Qwen3-1.7B & 0.0 / 0.0 \\
Qwen3-4B & 18.8 / 25.7 \\
Qwen3-8B & 0.0 / 0.2 \\
\end{longtable}
}

\subsection{Discussion}\label{discussion}

\textbf{Readability is not writability, and scale widens the gap.} The linear representation hypothesis
holds that concepts appear as directions \citep{park2023linear}; our readability axis measures exactly
that, and it climbs sharply with model size in both families. Yet the causal test --- does adding the
direction change behavior --- does not follow. Qwen3-4B (\emph{d}=2.07) and Qwen3-8B (\emph{d}=2.10) encode the
honest-vs-sycophantic distinction almost identically, but only the former is steerable. This is a
direct, cross-family instance of the point that ``representations that are readable from activations
may not be writable through them'' \citep{buchan2026dualstance}: a linear correlate of a concept need
not be a causal control knob, and scaling improves the correlate, not the control.

\textbf{Single-protocol steering evaluation is unsafe.} Had we evaluated only under A/B --- the common
default --- we would have concluded that CAA does nothing for sycophancy in any of six models, missing
the one genuine free-response effect. This is the steering-specific face of a general problem: models
are poor forced-choice selectors \citep{zheng2023mcqselectors}, and moving between multiple-choice and
open-style formats changes conclusions \citep{myrzakhan2024openllm, wang2025rightanswer}. Steering
claims should be reported under the protocol that matches the intended use (usually generation), with
a random-vector floor, and at a per-model operating point --- not a single imported layer/coefficient.

\textbf{What this does not claim.} It is not that steering never reduces sycophancy --- Qwen3-4B is a clear
positive --- but that the effect is rare, protocol-contingent, and not recovered by fair per-model
tuning where it is absent. We characterize \emph{where} and \emph{whether}, not the circuit-level \emph{why}; we
localize neither the components that make Qwen3-4B writable nor those that leave Qwen3-8B inert.
Identifying them --- e.g., via feature- or head-level analysis --- is the natural next step, with this
study's dissociation as its motivating observation.

\subsection{Limitations}\label{limitations}

\begin{itemize}
\tightlist
\item
  \textbf{Coefficient selection.} The coherence criterion (distinct-token ratio) is insensitive to
  fluent-but-derailed over-steering, so several models were run at the top of the coefficient range.
  Gemma-3-4B is null at a moderate coefficient (16) at its best layer, arguing against an
  over-steering artifact, but we did not repeat this check for every model.
\item
  \textbf{Power.} Baseline-sycophantic \emph{n} is low for Qwen3-8B under free-response (n=51), making that
  null the weakest-powered cell.
\item
  \textbf{Operating-point asymmetry.} Qwen3-4B, our sole positive, was validated at both a hand-set point
  and its outcome-blind swept point; other models were evaluated only at their swept points.
\item
  \textbf{A/B degeneracy.} The Qwen models prefer the sycophantic option on \textasciitilde100\% of A/B items at
  baseline and are never flipped, so the A/B null partly reflects an insensitive instrument rather
  than unchanged behavior --- which is itself part of the argument, but limits strong behavioral claims
  from A/B alone.
\item
  \textbf{Scope.} Six models are enough to show non-transfer and the dissociation but not to fit a scaling
  law; we study CAA only (not every steering method), a single judge model without human
  inter-rater agreement, and an item set that mixes MoralChoice-derived moral dilemmas and
  Perez-derived A/B pairs with custom-authored factual and opinion items. A full size ladder, ensemble judging
  with human calibration, and circuit-level localization are future work.
\end{itemize}

\subsection{Conclusion}\label{conclusion}

Across six models in two families, activation steering for sycophancy is measurement-fragile,
protocol-dependent, and largely non-transferring: it beats a random-vector floor in one of six models
under free-response and in none under A/B forced choice. The sycophancy direction becomes increasingly
\emph{readable} with scale while remaining, in most models, un-\emph{writable} --- the two axes are orthogonal,
most starkly in a pair of models with near-identical readability and opposite steerability. Reliable
steering evaluation therefore requires random-vector controls, protocol matched to intended use, and
per-model operating points; and the readability--writability gap is a load-bearing constraint on
representation-level control, not an artifact to tune away.

\subsection{Code and data availability}\label{code-and-data-availability}

All code, the item pool, and experiment outputs are released at
\textbf{https://github.com/barnalipics/readable-not-writable-sycophancy} (code MIT; data CC-BY-4.0 with
attribution for MoralChoice \citep{scherrer2023moralchoice} and the Anthropic model-written
evaluations \citep{perez2022discovering}). Every citation in this paper is verified programmatically
against the original source text (a fetch-and-substring-match script released with the code).

\subsection{Appendix}\label{appendix}

\subsubsection{A. Models and per-model operating points}\label{a.-models-and-per-model-operating-points}

All models are instruction-tuned and loaded with a single forward-hook addition at one decoder layer.
Readability is the maximum Cohen's \emph{d} (honest vs sycophantic activation projections, held-out split)
over the candidate layers; the operating point (layer, coef) is chosen outcome-blind (layer by
separation, coef by coherence frontier). Qwen loads in fp16, Gemma in bfloat16.

{\def\LTcaptype{none} % do not increment counter
\begin{longtable}[]{@{}
  >{\raggedright\arraybackslash}p{(\linewidth - 8\tabcolsep) * \real{0.2000}}
  >{\raggedright\arraybackslash}p{(\linewidth - 8\tabcolsep) * \real{0.2000}}
  >{\raggedright\arraybackslash}p{(\linewidth - 8\tabcolsep) * \real{0.2000}}
  >{\raggedright\arraybackslash}p{(\linewidth - 8\tabcolsep) * \real{0.2000}}
  >{\raggedright\arraybackslash}p{(\linewidth - 8\tabcolsep) * \real{0.2000}}@{}}
\toprule\noalign{}
\begin{minipage}[b]{\linewidth}\raggedright
Model
\end{minipage} & \begin{minipage}[b]{\linewidth}\raggedright
HF id
\end{minipage} & \begin{minipage}[b]{\linewidth}\raggedright
Layers
\end{minipage} & \begin{minipage}[b]{\linewidth}\raggedright
Op. point (L / coef)
\end{minipage} & \begin{minipage}[b]{\linewidth}\raggedright
Readability \emph{d}
\end{minipage} \\
\midrule\noalign{}
\endhead
\bottomrule\noalign{}
\endlastfoot
Qwen3-1.7B & Qwen/Qwen3-1.7B & 28 & 18 / 32 & 0.30 \\
Qwen3-4B & Qwen/Qwen3-4B-Instruct-2507 & --- & 23 / 32 (swept); 12 / 16 (primary) & 2.07 \\
Qwen3-8B & Qwen/Qwen3-8B & --- & 23 / 32 & 2.10 \\
Gemma-3-1B & unsloth/gemma-3-1b-it & 26 & 6 / 32 & −0.01 \\
Gemma-3-4B & unsloth/gemma-3-4b-it & 34 & 22 / 16 & 0.68 \\
Gemma-3-12B & unsloth/gemma-3-12b-it & 48 & 31 / 32 & 1.24 \\
\end{longtable}
}

\subsubsection{\texorpdfstring{B. Full per-model fix-rates (genuine-fix \%, with baseline-caved \emph{n})}{B. Full per-model fix-rates (genuine-fix \%, with baseline-caved n)}}\label{b.-full-per-model-fix-rates-genuine-fix-with-baseline-caved-n}

Fix-rate is computed only on items sycophantic at baseline in that protocol; ``rand'' is the
random same-norm vector floor. A real vector is writable when it exceeds rand by \textgreater5 points.

{\def\LTcaptype{none} % do not increment counter
\begin{longtable}[]{@{}
  >{\raggedright\arraybackslash}p{(\linewidth - 16\tabcolsep) * \real{0.1111}}
  >{\raggedright\arraybackslash}p{(\linewidth - 16\tabcolsep) * \real{0.1111}}
  >{\raggedright\arraybackslash}p{(\linewidth - 16\tabcolsep) * \real{0.1111}}
  >{\raggedright\arraybackslash}p{(\linewidth - 16\tabcolsep) * \real{0.1111}}
  >{\raggedright\arraybackslash}p{(\linewidth - 16\tabcolsep) * \real{0.1111}}
  >{\raggedright\arraybackslash}p{(\linewidth - 16\tabcolsep) * \real{0.1111}}
  >{\raggedright\arraybackslash}p{(\linewidth - 16\tabcolsep) * \real{0.1111}}
  >{\raggedright\arraybackslash}p{(\linewidth - 16\tabcolsep) * \real{0.1111}}
  >{\raggedright\arraybackslash}p{(\linewidth - 16\tabcolsep) * \real{0.1111}}@{}}
\toprule\noalign{}
\begin{minipage}[b]{\linewidth}\raggedright
Model
\end{minipage} & \begin{minipage}[b]{\linewidth}\raggedright
FF fac
\end{minipage} & \begin{minipage}[b]{\linewidth}\raggedright
FF rim
\end{minipage} & \begin{minipage}[b]{\linewidth}\raggedright
FF rand
\end{minipage} & \begin{minipage}[b]{\linewidth}\raggedright
FF \emph{n}
\end{minipage} & \begin{minipage}[b]{\linewidth}\raggedright
AB fac
\end{minipage} & \begin{minipage}[b]{\linewidth}\raggedright
AB rim
\end{minipage} & \begin{minipage}[b]{\linewidth}\raggedright
AB rand
\end{minipage} & \begin{minipage}[b]{\linewidth}\raggedright
AB \emph{n}
\end{minipage} \\
\midrule\noalign{}
\endhead
\bottomrule\noalign{}
\endlastfoot
Qwen3-1.7B & 9.2 & 8.2 & 13.3 & 98 & 0.0 & 0.0 & 0.0 & 463 \\
Qwen3-4B (swept) & 10.4 & \textbf{34.1} & 9.5 & 211 & 7.9 & 18.8 & 25.7 & 101 \\
Qwen3-4B (primary L12/c16) & \textbf{32.2} & \textbf{28.4} & 12.3 & 211 & 6.1 & 24.2 & 22.2 & 99 \\
Qwen3-8B & 9.8 & 11.8 & 13.7 & 51 & 0.0 & 0.0 & 0.2 & 463 \\
Gemma-3-1B & 4.8 & 7.3 & 9.5 & 357 & 8.3 & 8.3 & 7.6 & 132 \\
Gemma-3-4B & 2.2 & 1.9 & 4.8 & 314 & 0.0 & 0.0 & 1.2 & 82 \\
Gemma-3-12B & 4.6 & 3.5 & 3.5 & 284 & 0.0 & 1.6 & 0.0 & 62 \\
\end{longtable}
}

FF = free-response, AB = A/B forced choice. Only Qwen3-4B exceeds the random floor (free-response);
A/B is null for all models.

\subsubsection{C. Protocols, vectors, and scoring}\label{c.-protocols-vectors-and-scoring}

\textbf{Free-response (two-turn).} The reference answer is seeded as the assistant's first turn; the user
then asserts the contested claim (pushback); the model generates (greedy, 90 new tokens, per item).
An LLM judge (Gemini-2.5-flash) scores the reply on a reaffirmation-required rubric: a \emph{genuine fix}
requires the model to reaffirm the defensible position without evasion; contestable (opinion/moral)
items use a variant in which nuance counts as holding, not evasion. Baseline severity ≥ 2 marks an
item sycophantic at baseline.

\textbf{A/B forced choice.} The reference and contested claim are rendered as options (A)/(B), each item in
both orders; the next-token logit over the option letters gives the preference. An item is baseline-
sycophantic if it prefers the contested option; a fix is a flip to the reference preference. No judge.

\textbf{Vectors.} Factual (CAA from free-form contrast pairs), rimsky-style (CAA from A/B contrast pairs of
the model-written sycophancy set), and a random same-norm control; all unit-normalized and added at
the chosen layer with coefficient \emph{c}.

\subsubsection{D. Statistical details}\label{d.-statistical-details}

The protocol study is pre-registered. The primary test is a GEE (Binomial family, exchangeable working
correlation, clusters = item) for the protocol × vector interaction on genuine-fix; secondary tests
are per-protocol paired McNemar (rim vs fac) and bootstrap 95\% CIs on fix-rate differences. The
random same-norm vector is the validity floor in every cell. Reported cross-model comparisons are
descriptive fix-rate vs floor at each model's operating point.

\subsubsection{E. Reproducibility}\label{e.-reproducibility}

Generation is strictly per-item (batched generation corrupts the measurement). Gemma is loaded in
bfloat16 (fp16 yields NaN activations on our accelerators); Qwen in fp16. Large-model runs (Qwen3-8B,
Gemma-3-12B) used a single 32 GB GPU. Exact model ids, prompts, the rubric text, the shared item pool,
per-model selection logs, and all outputs are in the released repository (Appendix ``Code and data'').

\bibliographystyle{plainnat}
\bibliography{references_verified}
\end{document}
